\documentclass[10pt,a4paper]{amsart}
\usepackage{amsmath,amssymb,amsthm,graphicx}
% Retained minimal preamble for the Stacks Project chapter varieties.tex.
% Verbatim excerpt of the chapter preamble: Xy-pic setup (native preamble lines 14-18), the native theorem
% declarations with the mathematical macros (native preamble lines 154-217), and the project own no-op
% definitions of the reference/slogan/history display environments (native wrapper.tex lines 1-3) that the
% counted statement uses for its slogan. Copied mechanically; nothing added, removed or redefined.
% For commutative diagrams we use Xy-pic
\usepackage[all]{xy}

% We use 2cell for 2-commutative diagrams.
\xyoption{2cell}
% Theorem environments.
%
\theoremstyle{plain}
\newtheorem{theorem}[subsection]{Theorem}
\newtheorem{proposition}[subsection]{Proposition}
\newtheorem{lemma}[subsection]{Lemma}

\theoremstyle{definition}
\newtheorem{definition}[subsection]{Definition}
\newtheorem{example}[subsection]{Example}
\newtheorem{exercise}[subsection]{Exercise}
\newtheorem{situation}[subsection]{Situation}

\theoremstyle{remark}
\newtheorem{remark}[subsection]{Remark}
\newtheorem{remarks}[subsection]{Remarks}

\numberwithin{equation}{subsection}

% Macros
%
\def\lim{\mathop{\mathrm{lim}}\nolimits}
\def\colim{\mathop{\mathrm{colim}}\nolimits}
\def\Spec{\mathop{\mathrm{Spec}}}
\def\Hom{\mathop{\mathrm{Hom}}\nolimits}
\def\Ext{\mathop{\mathrm{Ext}}\nolimits}
\def\SheafHom{\mathop{\mathcal{H}\!\mathit{om}}\nolimits}
\def\SheafExt{\mathop{\mathcal{E}\!\mathit{xt}}\nolimits}
\def\Sch{\mathit{Sch}}
\def\Mor{\mathop{\mathrm{Mor}}\nolimits}
\def\Ob{\mathop{\mathrm{Ob}}\nolimits}
\def\Sh{\mathop{\mathit{Sh}}\nolimits}
\def\NL{\mathop{N\!L}\nolimits}
\def\CH{\mathop{\mathrm{CH}}\nolimits}
\def\proetale{{pro\text{-}\acute{e}tale}}
\def\etale{{\acute{e}tale}}
\def\QCoh{\mathit{QCoh}}
\def\Ker{\mathop{\mathrm{Ker}}}
\def\Im{\mathop{\mathrm{Im}}}
\def\Coker{\mathop{\mathrm{Coker}}}
\def\Coim{\mathop{\mathrm{Coim}}}

% Boxtimes
%
\DeclareMathSymbol{\boxtimes}{\mathbin}{AMSa}{"02}

%
% Macros for moduli stacks/spaces
%
\def\QCohstack{\mathcal{QC}\!\mathit{oh}}
\def\Cohstack{\mathcal{C}\!\mathit{oh}}
\def\Spacesstack{\mathcal{S}\!\mathit{paces}}
\def\Quotfunctor{\mathrm{Quot}}
\def\Hilbfunctor{\mathrm{Hilb}}
\def\Curvesstack{\mathcal{C}\!\mathit{urves}}
\def\Polarizedstack{\mathcal{P}\!\mathit{olarized}}
\def\Complexesstack{\mathcal{C}\!\mathit{omplexes}}
% \Pic is the operator that assigns to X its picard group, usage \Pic(X)
% \Picardstack_{X/B} denotes the Picard stack of X over B
% \Picardfunctor_{X/B} denotes the Picard functor of X over B
\def\Pic{\mathop{\mathrm{Pic}}\nolimits}
\def\Picardstack{\mathcal{P}\!\mathit{ic}}
\def\Picardfunctor{\mathrm{Pic}}
\def\Deformationcategory{\mathcal{D}\!\mathit{ef}}
\newenvironment{reference}{}{}
\newenvironment{slogan}{}{}
\newenvironment{history}{}{}
\usepackage{cleveref}
\crefname{originallemma}{original lemma}{original lemmas}
\Crefname{originallemma}{Original lemma}{Original lemmas}
\makeatletter
\newlabel{morphisms-lemma-algebraic-residue-field-extension-closed-point-fibre}{{01TE}{source}{}{sourcepointer-1}{}}
\newlabel{morphisms-lemma-algebraic-residue-field-extension-closed-point-fibre@cref}{{[originallemma][0][]01TE}{[1][0][]source}{}{}{}}
\makeatother
\crefname{originallemma}{original lemma}{original lemmas}
\Crefname{originallemma}{Original lemma}{Original lemmas}
\makeatletter
\newlabel{morphisms-lemma-base-change-quasi-finite}{{01TM}{source}{}{sourcepointer-2}{}}
\newlabel{morphisms-lemma-base-change-quasi-finite@cref}{{[originallemma][0][]01TM}{[1][0][]source}{}{}{}}
\makeatother
\crefname{originallemma}{original lemma}{original lemmas}
\Crefname{originallemma}{Original lemma}{Original lemmas}
\makeatletter
\newlabel{morphisms-lemma-dimension-formula}{{02JU}{source}{}{sourcepointer-3}{}}
\newlabel{morphisms-lemma-dimension-formula@cref}{{[originallemma][0][]02JU}{[1][0][]source}{}{}{}}
\makeatother
\crefname{originallemma}{original lemma}{original lemmas}
\Crefname{originallemma}{Original lemma}{Original lemmas}
\makeatletter
\newlabel{morphisms-lemma-finite-type-noetherian}{{01T6}{source}{}{sourcepointer-4}{}}
\newlabel{morphisms-lemma-finite-type-noetherian@cref}{{[originallemma][0][]01T6}{[1][0][]source}{}{}{}}
\makeatother
\crefname{originallemma}{original lemma}{original lemmas}
\Crefname{originallemma}{Original lemma}{Original lemmas}
\makeatletter
\newlabel{morphisms-lemma-quasi-finite-at-point-characterize}{{01TH}{source}{}{sourcepointer-5}{}}
\newlabel{morphisms-lemma-quasi-finite-at-point-characterize@cref}{{[originallemma][0][]01TH}{[1][0][]source}{}{}{}}
\makeatother
\crefname{originallemma}{original lemma}{original lemmas}
\Crefname{originallemma}{Original lemma}{Original lemmas}
\makeatletter
\newlabel{morphisms-lemma-quasi-finite-points-open}{{01TI}{source}{}{sourcepointer-6}{}}
\newlabel{morphisms-lemma-quasi-finite-points-open@cref}{{[originallemma][0][]01TI}{[1][0][]source}{}{}{}}
\makeatother
\crefname{originallemma}{original lemma}{original lemmas}
\Crefname{originallemma}{Original lemma}{Original lemmas}
\makeatletter
\newlabel{morphisms-lemma-residue-field-quasi-finite}{{01TG}{source}{}{sourcepointer-7}{}}
\newlabel{morphisms-lemma-residue-field-quasi-finite@cref}{{[originallemma][0][]01TG}{[1][0][]source}{}{}{}}
\makeatother
\crefname{originallemma}{original lemma}{original lemmas}
\Crefname{originallemma}{Original lemma}{Original lemmas}
\makeatletter
\newlabel{schemes-lemma-specialize-points}{{01J7}{source}{}{sourcepointer-8}{}}
\newlabel{schemes-lemma-specialize-points@cref}{{[originallemma][0][]01J7}{[1][0][]source}{}{}{}}
\makeatother
\title{Quasi-finiteness at points over a one-dimensional local ring from conditions at generic points}
\author{The Stacks Project Authors}
\date{stacks-project:varieties.tex:local-snapshot-20261011; GNU FDL 1.2 or later}
\begin{document}
\maketitle
\noindent\textit{Editorial note.} All delivered mathematical text is the author's own, copied from the pinned Stacks Project chapter file varieties.tex; no statement or proof step was rewritten, reordered or removed.\par\medskip
\noindent\textit{Editorial note.} The list of hypotheses in the statement ends with the author's own placeholder 'add more here'; that item is reproduced verbatim and no condition was added, removed or reordered.\par\medskip
\noindent\textit{Editorial note.} The counted spans carry 9 navigation stubs and no bibliography citation.\par\medskip
\noindent\textit{Editorial note.} The preamble is the chapter's retained minimal preamble (sources/stacks-varieties/retained-preamble.tex lines 1-77), a verbatim excerpt of the native preamble.\par\medskip
\noindent\textit{Editorial note.} The label lemma-quasi-finite-in-codim-1 is defined only in varieties.tex, so no cross-chapter restatement had to be weighed for this card.\par\medskip
\noindent\textit{Allowed standard prerequisites (editorial).}\begin{itemize}
\item Quasi-finite locus of a locally of finite type morphism is open. Morphisms, Lemma morphisms-lemma-quasi-finite-points-open, invoked at proof lines 3211-3212. Supplies that a dense set of points at which f is quasi-finite contains a dense open, hence all generic points of irreducible components.
\item Quasi-finiteness at a point versus finiteness of the residue field extension. Morphisms, Lemma morphisms-lemma-residue-field-quasi-finite and Lemma morphisms-lemma-quasi-finite-at-point-characterize, invoked at proof lines 3215-3216, 3235-3236 and 3263-3264. Supplies the translation between the hypotheses at generic points and the conclusion at fibre points.
\item Spec of a local ring as the set of points specializing to the point, and base change of quasi-finiteness. Schemes, Lemma schemes-lemma-specialize-points and Morphisms, Lemma morphisms-lemma-base-change-quasi-finite, invoked at proof lines 3221-3225. Supplies the reduction that lets Y be assumed to be the spectrum of a Noetherian local ring of dimension at most one.
\item Dimension formula for a morphism of schemes. Morphisms, Lemma morphisms-lemma-dimension-formula, invoked at proof lines 3246-3247. Supplies the inequality from which the vanishing of the transcendence degree is read off.
\end{itemize}\medskip
\ifdefined\hypertarget\hypertarget{sourcepointer-1}{}\fi
\noindent\textit{Source reference.} Original lemma 01TE: sources/stacks-chow/tags.txt, source lines 2371--2371. Pointer outside the counted span, invoked once for context; the load-bearing argument of the counted statement is the author's proof reproduced in full in the proof span.\par\medskip
\ifdefined\hypertarget\hypertarget{sourcepointer-2}{}\fi
\noindent\textit{Source reference.} Original lemma 01TM: sources/stacks-chow/tags.txt, source lines 2379--2379. Pointer outside the counted span, invoked once for context; the load-bearing argument of the counted statement is the author's proof reproduced in full in the proof span.\par\medskip
\ifdefined\hypertarget\hypertarget{sourcepointer-3}{}\fi
\noindent\textit{Source reference.} Original lemma 02JU: sources/stacks-chow/tags.txt, source lines 3349--3349. Pointer outside the counted span, invoked once for context; the load-bearing argument of the counted statement is the author's proof reproduced in full in the proof span.\par\medskip
\ifdefined\hypertarget\hypertarget{sourcepointer-4}{}\fi
\noindent\textit{Source reference.} Original lemma 01T6: sources/stacks-chow/tags.txt, source lines 2363--2363. Pointer outside the counted span, invoked once for context; the load-bearing argument of the counted statement is the author's proof reproduced in full in the proof span.\par\medskip
\ifdefined\hypertarget\hypertarget{sourcepointer-5}{}\fi
\noindent\textit{Source reference.} Original lemma 01TH: sources/stacks-chow/tags.txt, source lines 2374--2374. Pointer outside the counted span, invoked once for context; the load-bearing argument of the counted statement is the author's proof reproduced in full in the proof span.\par\medskip
\ifdefined\hypertarget\hypertarget{sourcepointer-6}{}\fi
\noindent\textit{Source reference.} Original lemma 01TI: sources/stacks-chow/tags.txt, source lines 2375--2375. Pointer outside the counted span, invoked once for context; the load-bearing argument of the counted statement is the author's proof reproduced in full in the proof span.\par\medskip
\ifdefined\hypertarget\hypertarget{sourcepointer-7}{}\fi
\noindent\textit{Source reference.} Original lemma 01TG: sources/stacks-chow/tags.txt, source lines 2373--2373. Pointer outside the counted span, invoked once for context; the load-bearing argument of the counted statement is the author's proof reproduced in full in the proof span.\par\medskip
\ifdefined\hypertarget\hypertarget{sourcepointer-8}{}\fi
\noindent\textit{Source reference.} Original lemma 01J7: sources/stacks-chow/tags.txt, source lines 2004--2004. Pointer outside the counted span, invoked once for context; the load-bearing argument of the counted statement is the author's proof reproduced in full in the proof span.\par\medskip
\begin{lemma}
\label{lemma-quasi-finite-in-codim-1}
Let $f : X \to Y$ be locally of finite type. Let $y \in Y$ be a point
such that $\mathcal{O}_{Y, y}$ is Noetherian of dimension $\leq 1$.
Assume in addition one of the following conditions is satisfied
\begin{enumerate}
\item for every generic point $\eta$ of an irreducible component
of $X$ the field extension $\kappa(\eta)/\kappa(f(\eta))$
is finite (or algebraic),
\item for every generic point $\eta$ of an irreducible component
of $X$ such that $f(\eta) \leadsto y$ the field extension
$\kappa(\eta)/\kappa(f(\eta))$ is finite (or algebraic),
\item $f$ is quasi-finite at every generic point of an
irreducible component of $X$,
\item $Y$ is locally Noetherian and $f$
is quasi-finite at a dense set of points of $X$,
\item add more here.
\end{enumerate}
Then $f$ is quasi-finite at every point of $X$ lying over $y$.
\end{lemma}

\begin{proof}
Condition (4) implies $X$ is locally Noetherian
(Morphisms, Lemma \ref{morphisms-lemma-finite-type-noetherian}).
The set of points at which morphism is quasi-finite is open
(Morphisms, Lemma \ref{morphisms-lemma-quasi-finite-points-open}).
A dense open of a locally Noetherian scheme contains all generic
point of irreducible components, hence (4) implies (3).
Condition (3) implies condition
(1) by Morphisms, Lemma \ref{morphisms-lemma-residue-field-quasi-finite}.
Condition (1) implies condition (2).
Thus it suffices to prove the lemma in case (2) holds.

\medskip\noindent
Assume (2) holds. Recall that $\Spec(\mathcal{O}_{Y, y})$
is the set of points of $Y$ specializing to $y$, see
Schemes, Lemma \ref{schemes-lemma-specialize-points}.
Combined with
Morphisms, Lemma \ref{morphisms-lemma-base-change-quasi-finite}
this shows we may replace $Y$ by $\Spec(\mathcal{O}_{Y, y})$.
Thus we may assume $Y = \Spec(B)$ where $B$ is a Noetherian
local ring of dimension $\leq 1$ and $y$ is the closed point.

\medskip\noindent
Let $X = \bigcup X_i$ be the irreducible components of $X$ viewed
as reduced closed subschemes. If we can show each fibre $X_{i, y}$
is a discrete space, then $X_y = \bigcup X_{i, y}$ is discrete as
well and we conclude that $X \to Y$ is quasi-finite at all points
of $X_y$ by Morphisms, Lemma
\ref{morphisms-lemma-quasi-finite-at-point-characterize}.
Thus we may assume $X$ is an integral scheme.

\medskip\noindent
If $X \to Y$ maps the generic point $\eta$ of $X$ to $y$, then $X$
is the spectrum of a finite extension of $\kappa(y)$ and the
result is true. Assume that $X$ maps $\eta$ to a point corresponding
to a minimal prime $\mathfrak q$ of $B$ different from $\mathfrak m_B$.
We obtain a factorization $X \to \Spec(B/\mathfrak q) \to \Spec(B)$.
Let $x \in X$ be a point lying over $y$.
By the dimension formula
(Morphisms, Lemma \ref{morphisms-lemma-dimension-formula})
we have
$$
\dim(\mathcal{O}_{X, x}) \leq \dim(B/\mathfrak q) +
\text{trdeg}_{\kappa(\mathfrak q)}(R(X)) - \text{trdeg}_{\kappa(y)} \kappa(x)
$$
We know that $\dim(B/\mathfrak q) = 1$, that the generic point of $X$
is not equal to $x$ and specializes to $x$ and that $R(X)$ is algebraic
over $\kappa(\mathfrak q)$. Thus we get
$$
1 \leq 1 - \text{trdeg}_{\kappa(y)} \kappa(x)
$$
Hence every point $x$ of $X_y$ is closed in $X_y$ by
Morphisms, Lemma
\ref{morphisms-lemma-algebraic-residue-field-extension-closed-point-fibre}
and hence $X \to Y$ is quasi-finite at every point $x$ of $X_y$ by
Morphisms, Lemma \ref{morphisms-lemma-quasi-finite-at-point-characterize}
(which also implies that $X_y$ is a discrete topological space).
\end{proof}

\end{document}
